% Chapter 1: Introduction
\chapter{서론}
\label{ch:introduction}

\section{연구 배경}
\label{sec:intro_background}

2019년 말 출현한 SARS-CoV-2는 전 세계적으로 7억 건 이상의 확진 사례를 발생시키며, 바이러스 유전체 분석의 중요성을 재확인시켰다~\cite{harvey2021tracking}.
바이러스는 복제 과정에서 지속적으로 변이를 축적하며, 특정 유전체 위치에 변이가 집중되는 현상을 hotspot이라 한다.
Hotspot은 전파력 증가, 면역 회피(immune evasion, 기존 면역 반응을 피하는 능력), 약물 저항성 등 바이러스의 주요 표현형 변화와 직접 관련되므로, 이를 정확히 검출하는 것은 백신 설계, 항바이러스제 개발, 실시간 유전체 감시의 핵심 과제이다~\cite{youn2025mutclust,nextstrain2018hadfield,rambaut2020pangolin}.

그러나 바이러스 변이 hotspot 검출 방법들의 성능을 체계적으로 비교·평가할 수 있는 벤치마크 프레임워크가 확립된 바 없다. 형식적으로, 바이러스 단백질 서열의 다중 서열 정렬(MSA)이 주어졌을 때, 검출 과제는 통계적으로 유의하게 높은 변이율을 보이는 유전체 위치의 부분집합을 식별하는 것이다. 그러나 표준화된 ground truth, 복합 평가 지표, 교차 병원체 검증 방법론이 존재하지 않으며, 이는 암 유전체학의 부분적 벤치마킹~\cite{bailey2018comprehensive}에도 불구하고 마찬가지이다.
본 연구자는 이전에 다중 오믹스 통합을 위한 체계적 벤치마크 프레임워크(MOSD~\cite{han2025mosd})와 적응적 hotspot 검출 알고리즘(MutClust~\cite{youn2025mutclust})을 개발하였으며, 정량적 평가 프레임워크 없이 MutClust를 개발한 경험이 본 연구의 직접적 동기가 되었다.

Nextstrain~\cite{nextstrain2018hadfield}이나 PANGO 계통 분류~\cite{rambaut2020pangolin}와 같은 현행 바이러스 유전체 플랫폼은 변이체 출현을 모니터링하지만, hotspot 검출 방법의 체계적 평가를 제공하지 않는다.
표준화된 벤치마크의 부재는 방법 개발자가 자신의 접근법을 객관적으로 비교할 수 없고, 최종 사용자가 특정 병원체에 어떤 방법을 적용할지 정보에 기반한 선택을 할 수 없음을 의미한다.


\section{연구 동기 및 문제 정의}
\label{sec:intro_motivation}

바이러스 변이 hotspot 검출의 체계적 평가가 부재함으로 인해, 다음과 같은 세 가지 구체적 문제가 발생한다.

첫째, \textbf{표준화된 평가 프레임워크의 부재}이다.
MutClust를 포함한 기존 hotspot 검출 방법들은 각자의 데이터셋과 평가 기준을 사용하여 성능을 보고하므로, 방법 간 공정한 비교가 불가능하다.
MutClust의 유일한 검증 수단인 bootstrap 검정은 검출된 클러스터의 통계적 유의성만을 확인하는 자기 참조적 방법으로, 해당 클러스터가 생물학적으로 의미 있는 기능적 영역에 대응하는지는 검증할 수 없다.

둘째, \textbf{독립적 ground truth에 기반한 검증의 부재}이다.
Hotspot 검출 결과를 객관적으로 평가하려면 알고리즘과 독립적인 정답 기준이 필요하나, 바이러스 유전체학에서 이러한 ground truth 체계가 정의된 사례가 없다.
수렴 진화(convergent evolution, 독립적 계통에서 동일 위치에 반복적으로 발생하는 변이) 위치, 정제 선택(purifying selection, 해로운 변이를 제거하는 자연선택) 위치, Deep Mutational Scanning(DMS, 단백질의 모든 가능한 단일 아미노산 변이의 기능적 영향을 실험적으로 측정하는 기술) 데이터 등 활용 가능한 독립적 증거원이 존재함에도 체계적으로 통합되지 않았다.

셋째, \textbf{최적 방법의 병원체 의존성에 대한 경험적 증거의 부재}이다.
기존 방법들은 단일 병원체에 맞추어 조정된 매개변수를 사용하며, 그 일반화 가능성이 통계적으로 확립되지 않았다.
병원체별 유전체 특성(길이, 변이율, 재조합 빈도 등)의 차이가 최적 검출 전략에 체계적 영향을 미치는지는 검증되지 않았다.
이것은 알고리즘 선택 문제(algorithm selection problem)~\cite{rice1976algorithm}의 구체적 사례이다: 특정 게놈 특성을 가진 병원체가 주어졌을 때, 어떤 탐지 방법을 적용해야 하는가? 체계적인 교차 병원체 평가 없이는 이 질문에 답할 수 없다.


\section{연구 목적}
\label{sec:intro_objectives}

위의 문제를 해결하기 위하여, 본 연구는 다음 두 가지 목표를 설정한다.

\begin{enumerate}[label=\textbf{목표 \arabic*.}]
  \item \textbf{바이러스 변이 hotspot 검출을 위한 체계적 벤치마크 프레임워크(MutBench)의 구축}: 3계층 ground truth(수렴 진화 기반 양성, 정제 선택 기반 음성, DMS 기능적 적합도), hotspot-score $= (\text{recall} + \text{precision} + \text{stability}) / 3$ (여기서 stability는 부분추출 반복 시 탐지 결과의 재현성을 측정한다) 복합 평가 지표, 2단계 실험 설계(1단계: SARS-CoV-2 심층 분석, 2단계: 9종 RNA 바이러스 $\times$ 9가지 스코어링 $\times$ 39개 탐지 방법, 총 3,159개 평가)를 포함하는 벤치마크 프레임워크를 구축한다.

  \item \textbf{병원체 적응적 방법 선택의 필요성에 대한 경험적 입증}: 대규모 벤치마크 결과에 대한 통계 분석(ANOVA, Friedman 검정, LOPO 교차 검증)을 통해, 단일 최적 방법이라는 전제가 병원체 전반에 걸쳐 성립하지 않음을 경험적으로 확립하고, PAHD 프레임워크의 개념 증명 프로토타입을 제시한다(제~\ref{ch:results}장 참조).
\end{enumerate}


\section{연구 기여}
\label{sec:intro_contributions}

본 연구의 학술적 기여는 다음과 같다.

\textbf{기여 1. MutBench 벤치마크 프레임워크.}
MutBench는 바이러스 변이 hotspot 검출을 위한 최초의 체계적 벤치마킹 프레임워크이다.
수렴 진화 위치를 양성 기준(Layer~A), 정제 선택 위치를 음성 기준(Layer~B), DMS 실험 데이터를 기능적 적합도 기준(Layer~C)으로 활용하는 3계층 ground truth 체계를 정의하였다.
정확도(recall, precision)와 신뢰성(stability)을 동시에 평가하는 hotspot-score 복합 지표를 도입하고, SARS-CoV-2 심층 분석과 9종 병원체 대규모 비교를 결합하는 2단계 실험 설계를 제시하였다.

\textbf{기여 2. 병원체 적응적 방법 선택 필요성의 통계적 입증.}
9종 RNA 바이러스에 대한 평가 결과, 9개 병원체 모두 상이한 최적 조합을 보이며(LOPO 교차 검증 0/9 일치), 방법 계열과 병원체 간 상호작용이 ANOVA에서 최대 분산 원천이었다.
Friedman 검정은 방법 간 성능 순위가 병원체에 따라 체계적으로 달라짐을 확인하였다.
이 결과는 단일 최적 방법이라는 전제가 성립하지 않음을 확립하며, PAHD 프레임워크의 개념 증명 프로토타입을 제시한다(상세 통계는 제~\ref{ch:results}장 참조).


\section{논문 구성}
\label{sec:intro_organization}

본 학위논문은 6개의 장으로 구성된다.

\textbf{제~\ref{ch:background}장, 관련 연구}에서는 바이러스 변이 hotspot 검출에 관한 기존 문헌을 조사하고, MutClust~\cite{youn2025mutclust}의 핵심 방법론 및 MutBench 개발의 동기가 된 한계점을 요약한다.

\textbf{제~\ref{ch:mutbench}장, 자료 및 방법}에서는 본 학위논문의 핵심 기여인 MutBench 벤치마크 프레임워크의 자료와 방법론을 상세히 기술한다.
제3.1절은 자료(대상 병원체, ground truth 체계, DMS 데이터)를 기술한다. 제3.2--3.4절은 방법을 제시한다: 제3.2절은 MutClust 변형 및 hotspot-score 지표를 포함하는 벤치마크 프레임워크 설계, 제3.3절은 스코어링 공식과 탐지 방법 계열, 제3.4절은 평가 지표 및 통계 분석 설계를 다룬다.

\textbf{제~\ref{ch:results}장, 실험 결과}에서는 실험 결과를 네 개 절로 구분하여 보고한다: 단일 병원체 벤치마크(합성 데이터 비교, 다중 ground truth 평가, 실제 GISAID 검증, SARS-CoV-2에 대한 외부 방법 비교), 견고성 검증(통계적 검증, 민감도 분석, 시간적 일반화, indel 제외), 교차 병원체 검증(H3N2 검증, 3D 구조 분석, 계통발생학적 시뮬레이션, ESM-2 교차 패러다임 검증), 9-병원체 대규모 벤치마크(ANOVA, Friedman 검정, LOPO 교차 검증, PAHD 개념 증명 프로토타입).

\textbf{제~\ref{ch:discussion}장, 논의}에서는 실험 결과의 종합적 해석, MutBench의 방법론적 함의, 한계점, 그리고 PAHD 개념적 프레임워크의 향후 방향을 논의한다.

\textbf{제~\ref{ch:conclusion}장, 결론}에서는 연구의 기여를 요약하고, 한계점을 분석하며, 향후 연구 방향을 제시한다.
